\documentclass[10pt]{article}
\usepackage{titling}

\title{Bonds}
\date{DATE PLACEHOLDER}
\author{Ingyu Bahng}

\usepackage[letterpaper, margin=1in]{geometry}
\usepackage{amsmath} % better math functions
\usepackage{amssymb} % better math symbols
\usepackage{babel}[english] % change rules to english
\usepackage{lipsum} % allows for lorem ipsum text generation
\usepackage{xr-hyper} % allows cross referencing labels from external LaTeX documents 
\usepackage{hyperref} % for hyperlinking text
\usepackage{fancyvrb} % for better verbatim environments
\usepackage{newverbs} % allows for inline typewriter font via \texttt{}
\usepackage{fancyvrb} % also code font blocks but has more capabilities than texttt
\usepackage{footnote} % allows footnotes inside tabulars
\usepackage{dcolumn} % allows for decimal alignment in tables
\usepackage{tabularx}

\usepackage{fontspec}
\setmainfont{Gentium Plus} % allowing greek letters in text

\usepackage{unicode-math}
\setmathfont{XITS Math} % allowing greek letters in math



% tcolorbox package
\usepackage{tcolorbox}
\tcbuselibrary{breakable}

\tcbset{ % this is the default design of tcolorboxes
  colframe = black,
  colback  = black!1,
  coltitle = black,
  colbacktitle = black!15,
  breakable, 
  arc=0mm,
  boxrule=0.5pt,
  toptitle=3pt,
  bottomtitle=3pt,
  left=8pt,
  right=8pt,
  top=6pt,
  bottom=6pt,
  before skip=12pt,
  after skip=12pt,
  bottomrule at break=-1pt,
  toprule at break=-1pt,
  fonttitle=\bfseries,
}

% Custom Environments
\newtcolorbox[auto counter, number within=section]{definition}[1][]
{
    title = \textbf{Definition \thetcbcounter: ~#1}
}

\newtcolorbox[auto counter, number within=section]{core}[1][]
{
    title = \textbf{Core Concept \thetcbcounter: ~#1}
}


\usepackage{fancyhdr}
\usepackage[skip=0.8em, indent=0em]{parskip}

\pagestyle{fancy} % this section generates the lines and text at the top and bottom of each page
\fancyhead[L]{\thetitle}
\fancyhead[C]{\theauthor}
\fancyhead[R]{\thedate}
\fancyfoot[C]{\thepage}
\renewcommand{\footrulewidth}{0.3pt}
\renewcommand{\thispagestyle}[1]{}

\renewcommand{\arraystretch}{1.5} % table vertical padding
\setlength{\tabcolsep}{10pt} % table horizontal padding

\usepackage{enumitem} % better control over enumerate and itemize

% List customization
\setlist[itemize]{%
  label=\bullet,
  itemsep=2pt,
  leftmargin=1.5em,
}

\setlist[enumerate]{%
  label=(\alph*),
  itemsep=2pt,
  leftmargin=1.5em,
}



\usepackage{pgfplots} % for plot creation
\pgfplotsset{compat=1.18} % setting the version used for pgfplots
\usepackage{float} % for better figure placement using [h]
\usepackage{graphicx} % for importing figures into the tex file
\usepackage{subcaption} % allows for multiple subfigures with individual captions wihtin one figure
\usepackage{xparse} % allows for more than 9 parameters in commands
\usepackage{adjustbox} % for valign option
\captionsetup{font=small} % setting the caption fontsize to small always

\usepackage{silence}
\WarningsOff[float]  % suppress float package warnings

\newcommand{\myfig}[4]{%
  \begin{figure}[H]
    \centering
    \adjustbox{valign=c}{\includegraphics[width=#1\textwidth]{#2}}
    \caption{#3}
    \label{#4}
  \end{figure}
}

\newcommand{\mymultifig}[8]{
  \begin{figure}[H]%
    \centering%
    \begin{subfigure}{#1\textwidth}%
      \centering%
      \includegraphics[width=\linewidth]{#2}%
    \end{subfigure}%
    \hfill%
    \begin{subfigure}{#3\textwidth}%
      \centering%
      \includegraphics[width=\linewidth]{#4}%
    \end{subfigure}%
    \hfill%
    \begin{subfigure}{#5\textwidth}%
      \centering%
      \includegraphics[width=\linewidth]{#6}%
    \end{subfigure}%
    \caption{#7}%
    \label{#8}%
  \end{figure}%
}

\usepackage{newverbs} % allows for inline typewriter font via \texttt{}
\usepackage{fancyvrb} % also code font blocks but has more capabilities than texttt
\usepackage{listings} % allows for codeblocks - config below

%BELOW ARE CODE BLOCK DESIGNS (Light Theme)
\definecolor{gray}{HTML}{707070}
\definecolor{lightgray}{HTML}{333333}
\definecolor{codebg}{HTML}{F5F5F5}
\definecolor{keyword}{HTML}{0000FF}
\definecolor{string}{HTML}{A31515}
\definecolor{number}{HTML}{098658}
\definecolor{function}{HTML}{795E26}
\definecolor{comment}{HTML}{008000}
\definecolor{classname}{HTML}{267F99}

% default options for listings (for code)
\lstset{
  frame=ltbr,
  language=python,
  aboveskip=3mm,
  belowskip=3mm,
  showstringspaces=false,
  columns=fullflexible,
  keepspaces=true,
  basicstyle=\fontsize{9}{10}\ttfamily\color{lightgray},
  numbers=left,
  firstnumber=1,
  numberstyle=\fontsize{7}{10}\ttfamily\color{gray},
  keywordstyle=\color{keyword},
  commentstyle=\color{comment},
  stringstyle=\color{string},
  backgroundcolor=\color{codebg},
  breaklines=true,
  breakatwhitespace=true,
  tabsize=2,
  xleftmargin=3em,
  numbersep=1em,
  framexleftmargin=2.5em,
  framesep=6pt,
  stepnumber=1,
}

\hfuzz=5.002pt % ignore overfull hbox badness warnings below this limit



\usepackage{chemfig}
\usepackage[version=4]{mhchem}

\newcommand{\bce}[1]{
  \hspace{7mm}{#1}
}

\newcommand{\sno}[0]{
  $\mathrm{S}_{\mathrm{N}}1$
}

\newcommand{\snt}[0]{
  $\mathrm{S}_{\mathrm{N}}2$
}

\newcommand{\reactionfig}[4]{
    \begin{figure}[H]
    \centering
    \scalebox{#4}{
    \schemestart
    #1
    \schemestop
    }
    \caption{#2}
    \label{#3}
    \end{figure}
}




\begin{document}

\input{../../presets/_general/startup}

\section{Bonds and Intermolecular Forces}

  \subsection{Bond Types}

  Properties of matter depend on the types of bonds and intermolecular forces (IMF) present. \textbf{Covalent bonds} involve the sharing of electrons between atoms. \textbf{Ionic bonds} result from the electrostatic attraction between oppositely charged ions. \textbf{Metallic bonds} are formed by the sharing of valence electrons. \textbf{IMFs}, such as hydrogen bonding, dipole-dipole interactions, and London dispersion forces, influence properties like boiling point, melting point, and solubility.

  1. \textbf{Molecular Covalent} (e.g. \ce{H2}, \ce{Cl2}, \ce{H2O}, \ce{CH4}) -- \textit{soft liquids}

    \myfig{0.5}{bonds_figures/molecular_covalent}{caption}{fig:fig1}

  2. \textbf{Network Covalent} (e.g. Carbon, Hydrocarbon, Sulfur) -- \textit{hard solids, non-brittle, non-conductive}

    \myfig{0.6}{bonds_figures/network_covalent}{caption}{fig:fig2}
  
  3. \textbf{Ionic} (e.g. \ce{NaCl}, \ce{MgCl2}, \ce{CuSO4}) -- \textit{anions(-) \& cations(+) are held together by electrostatic forces}

    \myfig{0.5}{bonds_figures/ionic}{caption}{fig:fig3}

  4. \textbf{Metallic} (e.g. \ce{Li}, \ce{Na}, \ce{Ca}, \ce{Al}) -- \textit{malleable, bendable, conductor}

  Metals consist of positively charged ions surrounded by a "sea of delocalized \ce{e^-}s", allowing for high electrical and thermal conductivity, malleability, and ductility.

    \myfig{0.25}{bonds_figures/metallic}{caption}{fig:fig4}
    \myfig{0.25}{bonds_figures/metallic}{caption}{fig:fig5}
    \myfig{0.25}{bonds_figures/metallic}{caption}{fig:fig6}

  \subsection{VSEPR and Molecular Geometry}
  
  The shapes of molecules can be predicted using the \textbf{VSEPR (Valence Shell Electron Pair Repulsion)} theory, which states that electron pairs around a central atom will arrange themselves to minimize repulsion.
  
    \myfig{0.9}{bonds_figures/VESPR_table}{caption}{fig:fig7}
  
  \ce{CH4} $\rightarrow$ no lone pairs, 4 \ce{e^-} dense areas $\rightarrow$ \textbf{tetrahedral} \\
  \ce{H2O} $\rightarrow$ 2 lone pairs, 4 \ce{e^-} dense areas $\rightarrow$ \textbf{bent} \\
  \ce{NH3} $\rightarrow$ 1 lone pair, 4 \ce{e^-} dense areas$\rightarrow$ \textbf{pyramidal}.

  \begin{definition}
  It is important to note here that \textbf{lone pairs take up more space than bonding pairs}, which can affect bond angles. For example, in \ce{NH3}, the bond angle is 107° instead of the ideal tetrahedral angle of 109.5° due to the presence of a lone pair on nitrogen. Similarly, in \ce{H2O}, the bond angle is 104.5° because of the two lone pairs on oxygen.
  \end{definition}

  \subsection{Bond Polarity and Electronegativity}

  The bond polarity of a molecule is determined by the difference in \textbf{electronegativity (EN)} between the bonded atoms. The strength of the bonds from weakest to strongest can be categorized as follows:

  1. \textbf{Nonpolar Covalent}: EN difference $< 0.5$ -- (\textit{e.g. \ce{H2}, \ce{Cl2}, \ce{F2}})

  2. \textbf{Polar Covalent}: EN difference between $0.5$ and $1.7$ --(\textit{e.g. \ce{H2O}, \ce{NH3}, \ce{HCl}})
  
  3. \textbf{Ionic}: EN difference $> 1.7$ -- (\textit{e.g. \ce{NaCl}, \ce{MgO}, \ce{CaF2}})

  The EN values of elements are fundamentally determined by the laws of electromagnetic forces, which depend on the \textbf{strength of opposing charges} (\ce{Q_1} \& \ce{Q_2}) and the \textbf{distance between them} (r). This is described by Coulomb's Law:
  
  \[
  F = k \frac{Q_1 Q_2}{r^2}
  \]
  
  \textbf{As you move across the periodic table from left to right ($\rightarrow$), EN increases} due to the increasing number of \ce{p^+} attracting the same number of \ce{e^-} shells. \textbf{As you move downwards ($\downarrow$), EN decreases} because the atomic radius increases, reducing the effective nuclear charge (\ce{Z_e_f_f}) felt by the valence electrons.

  The \ce{Z_{eff}} can be approximated using the formula
  
  \[
  Z_eff = Z - S
  \]
  
  where \ce{Z} is the actual nuclear charge (number of \ce{p^+}) and S is the number of \ce{e^-} within the valence shell (\textit{this excludes the valence \ce{e^-}s}).

  \begin{definition}
  Keep in mind that \textbf{the actual definition of \ce{Z_{eff}} is more nuanced} than the formula above. The shielding effect of inner electrons is not uniform, and different orbitals (s, p, d, f) have varying effects on \ce{Z_{eff}}. For example, s-electrons are more effective at shielding than p-electrons. Additionally, electron-electron repulsions within the same shell can also affect the effective nuclear charge experienced by valence electrons.
  \end{definition}

  \subsection{Intermolecular Forces}

  Before discussing intermolecular forces (IMFs), it is important to understand that although an individual bond may be polar, the \textbf{overall molecule can be nonpolar if the molecular geometry allows for the dipole moments to cancel out}. For example, \ce{CO2} has polar bonds, but the linear geometry results in a nonpolar molecule.
  
    \myfig{0.3}{bonds_figures/canceled_dipole}{caption}{fig:fig8}

  Thus, when evaluating molecular polarity, both bond polarity and molecular geometry must be considered.

  IMFs are the forces of attraction/repulsion \textbf{between} molecules. They are generally weaker than covalent or ionic bonds but play a crucial role in determining the physical properties of substances. The main types of IMFs include:

  1. \textbf{London Dispersion Forces (LDFs)} -- present in all molecules, but dominant in nonpolar molecules. They arise from temporary dipoles created by the movement of \ce{e^-}. The strength of LDFs increases with the size and shape of the molecule (\textit{more surface area = stronger LDFs}).
  
    \myfig{0.4}{bonds_figures/LDF}{caption}{fig:fig9}

  \begin{definition}
    The boiling point (BP) of alkanes (\ce{C_xH_y}) increases with increasing molar mass due to stronger LDFs. Similarly, a larger surface area in molecules like long-chain hydrocarbons leads to higher BPs compared to their more compact isomers due to an increased area for LDF interactions.

    \myfig{0.4}{bonds_figures/alkane_ldf}{caption}{fig:fig10}
  \end{definition}

  2. \textbf{Dipole-induced Dipole \& Ion-induced Dipole} -- occurs when a polar molecule or ion induces a dipole in a nonpolar molecule by distorting its electron cloud.

    \myfig{0.3}{bonds_figures/induced_dipole}{caption}{fig:fig11}

  3. \textbf{Dipole-Dipole Interactions} -- occur between polar molecules where the (+) end of one molecule is attracted to the (-) end of another.

    \myfig{0.45}{bonds_figures/dipole_dipole}{caption}{fig:fig12}

  4. \textbf{Hydrogen Bonding} -- a special case of dipole-dipole interaction that occurs when H is bonded to highly electronegative atoms like N, O, or F.

    \myfig{0.25}{bonds_figures/h_bonds}{caption}{fig:fig13}

  4. \textbf{Ion-Dipole Interactions} -- occur between ionic compounds and polar molecules, such as when salt (\ce{NaCl}) dissolves in water.

    \myfig{0.25}{bonds_figures/ion_dipole}{caption}{fig:fig14}

  \begin{definition}
  SIDE NOTE: \textbf{Formal Charge (FC)} is a useful concept for determining the most stable Lewis structure for a molecule. It is calculated using the formula

  \[
  FC = V - (L + \frac{B}{2})
  \]

  V $\rightarrow$ \# of valence \ce{e^-} \\
  L $\rightarrow$ \# of lone pairs \\
  B $\rightarrow$ \# of bonding \ce{e^-} \\
  
  The \textbf{most stable Lewis structure typically has the smallest formal charges} on atoms, with negative formal charges residing on the more electronegative atoms. \\
  
  \textit{For example, in \ce{NO3^-}, N has a formal charge of +1, while one of the O atoms has a formal charge of -1, and the other two O atoms have a formal charge of 0. This distribution of formal charges helps to explain the resonance structures and overall stability of the nitrate ion.}
  \end{definition}

\end{document}

